\section{Introduction}
  \label{sec:intro}

  \subsection{The stratigraphic problem}
    \label{sec:intro:problem}

    The Bounded Admissibility framework~\cite{Cosmochrony} associates to each
    relaxation rank $n$ an energy scale $E_n = E_{\mathrm{P}}/n$ and a spectral
    admissibility condition: a mode of eigenvalue $\lambda$ is projectable at rank
    $n$ if and only if $\lambda \leq \Lambda_{\mathrm{proj}}(n)$, where
    $\Lambda_{\mathrm{proj}}(n)$ is the projective resolution at that stage of the
    cascade.
    The \emph{stabilisation rank} of a mode is the first $n$ at which it ceases
    to be admissible, and the corresponding stabilisation energy is taken as a
    proxy for the physical mass of the associated particle sector.

    For this picture to produce a viable mass hierarchy, three conditions must hold
    simultaneously.
    First, the projective resolution must decrease with rank, so that earlier
    stabilisations correspond to higher masses.
    Second, the set of stabilisation ranks must be discrete, so that the spectrum
    of effective masses is not a continuum.
    Third, the inter-level separations must be large enough to account for the
    observed mass ratios between generations, which span five to seventeen orders
    of magnitude from the electron to the top quark.

    The first two conditions were addressed in~\cite{SpectralStratigraphy}.
    The third -- which requires a quantitative law for $\Lambda_{\mathrm{proj}}(n)$
    -- was left as an open problem there and is the subject of the present paper.

  \subsection{What SpectralStratigraphy established}
    \label{sec:intro:stratigraphy}

    The companion paper~\cite{SpectralStratigraphy} established two results that
    together constrain the form of $\Lambda_{\mathrm{proj}}(n)$ from above and
    below.

    \medskip
    \noindent\textbf{No-go for scalar calibration (Proposition~2
of~\cite{SpectralStratigraphy}).}
    A threshold of the form $E^*(\lambda) = E_0 \lambda^\beta$ cannot produce a
    discrete stratigraphic profile for any choice of parameters $(\alpha, \beta)$.
    This rules out any explanation of the discreteness of the generation spectrum
    by re-scaling or re-calibrating the energy-eigenvalue relation.
    The discretisation must therefore originate in the internal structure of the
    projective substrate.

    \medskip
    \noindent\textbf{Three-level structure from ADE binary groups
  (Proposition~3 of~\cite{SpectralStratigraphy}).}
    For four specific pairs $(G, S)$ drawn from the binary tetrahedral, octahedral,
    and icosahedral groups -- $2T$ with order-3 generators, $2O$ with order-3
    generators, and $2I$ with order-4 or order-5 generators -- the Cayley graph
    Laplacian has exactly three distinct non-zero eigenvalues.
    Each eigenvalue level corresponds to a block of irreducible representations.
    Consequently, the stratigraphy defined by the Born--Infeld admissibility
    envelope $A^{\max}(\lambda) \propto 1/\sqrt{\lambda}$ produces exactly three
    stratigraphic classes, ordered by decreasing admissibility window.

    These two results together establish the following separation of roles, which
    is the central logical outcome of~\cite{SpectralStratigraphy}:
    \begin{quote}
      \emph{Group topology determines the number of stratigraphic levels.
      Projective dynamics determines the inter-level separations.}
    \end{quote}

  \subsection{What remained open}
    \label{sec:intro:open}

    Despite the clean three-level structure, the numerical inter-level separations
    produced by the ADE examples are small: between $0.08$ and $0.18$ log-decades
    in cascade depth $\log n$, as reported in Table~2 of~\cite{SpectralStratigraphy}.
    This is far too small to account for mass ratios of order $10^5$--$10^{17}$.

    The gap is not a defect of the group-theoretic mechanism.
    It reflects the fact that the spectral ratios $\lambda_\ell / \lambda_h$ for
    the ADE Cayley graphs lie in the range $[1.5, 2.0]$: the eigenvalue levels are
    spectrally close.
    Amplifying these small spectral differences into large mass ratios requires a
    rapidly growing function $\Lambda_{\mathrm{proj}}(n)$.

    The missing ingredient is therefore precisely the \emph{growth law} of
    $\Lambda_{\mathrm{proj}}(n)$ along the relaxation trajectory.
    This law must be derived from within the framework, not imposed as an external
    ansatz: the no-go result rules out scalar power-law calibrations, and any
    other ad hoc choice would undermine the predictive content of the stratigraphy.

  \subsection{Aim and approach of the present paper}
    \label{sec:intro:aim}

    This paper proposes a derivation of $\Lambda_{\mathrm{proj}}(n)$ grounded in
    the global coherence properties of the relaxation graph $G(n)$.

    The key observation is that $\Lambda_{\mathrm{proj}}(n)$ is entirely controlled
    by the global projective flux bound $c_\chi(n)$:
    \[
      \Lambda_{\mathrm{proj}}(n) = \frac{c_\chi(n)^2}{A_{\min}^2}.
    \]
    The dynamical question reduces to: \emph{how does $c_\chi(n)$ depend on the
relaxation rank?}
    We argue that $c_\chi(n)$ is bounded by the isoperimetric capacity of $G(n)$,
    measured by the Cheeger constant $h(G(n))$, which is in turn controlled by the
    algebraic connectivity $\lambda_2(n)$ via the discrete Cheeger inequality.

    The logical chain of the derivation is:
    \[
      \text{global coherence}
      \;\longrightarrow\; c_\chi(n)
      \;\longrightarrow\; h\bigl(G(n)\bigr)
      \;\longrightarrow\; \lambda_2(n)
      \;\longrightarrow\; \Lambda_{\mathrm{proj}}(n).
    \]
    Two structural hypotheses are required: a \emph{projective coherence closure}
    (Hypothesis~\ref{hyp:closure}) and a \emph{spectral-isoperimetric saturation}
    condition (Hypothesis~\ref{hyp:saturation}).
    Both are stated explicitly and treated as independent structural postulates.

    The main results of the paper are as follows.
    \begin{itemize}
      \item \textbf{Proposition~\ref{prop:Lambda-h}}:
      $\Lambda_{\mathrm{proj}}(n) \asymp h(G(n))^2$ under Hypothesis~\ref{hyp:closure}
      alone.
      \item \textbf{Corollary~\ref{cor:cheeger-enclosure}}:
      $\lambda_2(n)^2 \lesssim \Lambda_{\mathrm{proj}}(n) \lesssim \lambda_2(n)$
      via the Cheeger inequality, with no additional hypotheses.
      \item \textbf{Corollary~\ref{cor:linear-law}}:
      $\Lambda_{\mathrm{proj}}(n) \asymp \lambda_2(n)$ under both hypotheses.
    \end{itemize}
    The programmatic consequence is that the problem of inter-generation mass ratios
    reduces to the problem of computing $\lambda_2(n)$ along the relaxation
    trajectory.

  \subsection{Structure of the paper}
    \label{sec:intro:structure}

    Section~\ref{sec:framework} establishes the formal framework.
    Section~\ref{sec:model} introduces the LPS graph model and identifies the
    relaxation index.
    Section~\ref{sec:spectral} analyses the spectral behaviour of LPS graphs and
    derives the growth law for $N(\lambda; n)$.
    Section~\ref{sec:amplification} evaluates the predicted mass ratios numerically.
    Section~\ref{sec:discussion} discusses limitations and open directions.
    Section~\ref{sec:conclusion} concludes.
